\chapter{The Open-Source MLOps Toolchain}
\label{ch:toolchain}

In order to construct a production-grade machine learning platform without incurring prohibitive enterprise cloud licensing costs, our architecture integrates a curated selection of open-source components. However, merely assembling disparate tools is insufficient. True operational resilience requires a deep, mechanical interlocking between data versioning, experiment telemetry, and hyperparameter optimisation.

\section{Content-Addressable Storage and Cryptographic Versioning}
\label{sec:dvc}

Standard version control systems, such as Git, are engineered to track line-by-line code modifications within directed acyclic graphs of text files. Applying this paradigm to binary biological datasets is architecturally catastrophic. Committing a 50-gigabyte archive of high-resolution images forces the repository database to duplicate entire assets upon the slightest modification. This inevitably leads to memory exhaustion and repository collapse.

The system safely circumvents this limitation by utilising Data Version Control (DVC), which implements content-addressable storage. Rather than tracking the physical image matrices, DVC computes a unique SHA-256 cryptographic hash based strictly on the pixel content of the target directory. The heavy binary payloads are routed directly to a remote storage bucket, while a lightweight pointer file (for example, \texttt{dataset.dvc}) containing only the resultant hash and file metadata is generated.

This pointer file is subsequently committed to Git. When a deployment runner or data scientist executes a checkout command for a specific historical commit, the DVC client reads the hash from the pointer file, resolves the asset location in the remote cache, and constructs near-instantaneous filesystem reflinks (or hardlinks) into the local workspace. This dual-layered approach reliably decouples the heavy binary storage from the codebase, whilst guaranteeing mathematically verifiable, bit-for-bit data reproducibility.

\section{The Unified Management Plane}
\label{sec:dagshub}

Operating a distributed MLOps toolchain traditionally requires provisioning disparate infrastructure: an AWS S3 bucket for data object storage, a PostgreSQL database for experiment telemetry, and a complex matrix of Identity and Access Management (IAM) policies in order to broker communication between them.

In order to eliminate this DevOps overhead, the architecture standardises on DagsHub as the central management plane. By mirroring the primary GitHub repository, DagsHub acts as an automated nexus that instantly provisions a secure, S3-compatible remote storage backend for DVC and a managed MLflow Tracking Server. This integration effectively collapses the required infrastructure footprint, allowing self-hosted execution runners and local development environments to authenticate and stream both multi-gigabyte datasets and real-time loss metrics through a unified API token gateway.

\section{Bayesian Hyperparameter Optimisation and Telemetry}
\label{sec:hyperparameter_opt}

Model convergence, latency minimisation, and epistemic reliability depend fundamentally on tuning discrete and continuous training hyperparameters. Rather than viewing hyperparameter configuration and experiment tracking as disjoint manual procedures, the architecture integrates Bayesian optimisation algorithms directly with distributed metric logging pipelines in order to form a self-contained optimisation loop.

\subsection{Intelligent Search Mechanics}

Neural network architectures require the optimisation of external hyperparameters, such as learning rates, weight decay penalties, and conformal confidence thresholds. Traditional Grid Search algorithms evaluate these spaces exhaustively, incurring a computational cost of $\mathcal{O}(S^D)$, where $S$ represents the number of discrete steps evaluated across $D$ hyperparameter dimensions. On a single-node bare-metal architecture, this geometric explosion renders hyperparameter tuning computationally unviable.

The platform replaces brute-force evaluation with Optuna, a framework driven by the Tree-structured Parzen Estimator (TPE). TPE abandons blind search grids in favour of Bayesian probability. Instead of modelling the objective function directly, it models the parameter probability distributions conditioned on historical trial performance. Formally, for a given hyperparameter configuration $\theta$ and a target loss metric $y$ (specifically configured to optimise the PR-AUC and AUPIMO benchmarks formally defined in Chapter 7), TPE models the conditional probability $P(\theta \mid y)$ as follows:
\[
P(\theta \mid y) = 
\begin{cases} 
    \ell(\theta) & \text{if } y < y^* \\ 
    g(\theta) & \text{if } y \ge y^* 
\end{cases}
\]
In this formulation:
\begin{itemize}
    \item $\theta$ is the specific hyperparameter configuration (for example, learning rate) being evaluated.
    \item $y$ is the empirical loss achieved by executing the network with configuration $\theta$.
    \item $y^*$ is a dynamically computed scalar quantile threshold separating high-performing trials from poor-performing trials.
    \item $\ell(\theta)$ is the probability density function formed by the parameters of successful past trials ($y < y^*$).
    \item $g(\theta)$ is the probability density function formed by the parameters of unsuccessful past trials ($y \ge y^*$).
\end{itemize}

The algorithm iteratively selects new configurations that maximise the ratio $\ell(\theta) / g(\theta)$. This guides the GPU to sample from regions mathematically correlated with low loss, whilst preserving enough variance to escape local minima.

In order to further conserve hardware cycles, Optuna executes the Asynchronous Successive Halving Algorithm (ASHA). If a trial's validation metrics fall into the bottom quartile during its initial training epochs, ASHA aggressively prunes (terminates) the run. This immediately frees the Tensor Cores to evaluate a more promising branch of the probability distribution.

\subsection{The MLflow Interlock Architecture}

While Optuna efficiently orchestrates the search mathematics locally, it inherently lacks the infrastructure to permanently archive model provenance across a distributed engineering team. MLflow bridges this gap, operating as the continuous flight data recorder for our training pipeline.

The two systems are strictly coupled via an \texttt{MLflowCallback} hook. As the pipeline executes, the overarching Optuna \textit{Study} is registered as a root experiment in MLflow. Every individual parameter combination evaluated by Optuna is recorded as a nested child run. During these nested executions, MLflow streams the step-by-step metrics (for example, training loss, top-1 accuracy) and the sampled parameters (for example, $\theta$) directly to the DagsHub dashboard. Upon completion, physical artefacts -- such as the serialised ONNX graphs and validation histograms -- are uploaded to the artefact registry. Programmatic triggers then evaluate the Pareto frontier of all logged trials, automatically promoting the optimal model checkpoint to the MLflow Model Registry, staging it for formal CI/CD deployment testing.
